\section{Gestion de Projet}

\subsection{Déroulement et organisation du projet}
Pour nous coordonner sur l’avancement du PSC nous nous sommes réunis tous les mercredis après-midi durant les créneaux réservés où nous pouvions alors échanger avec Clément Liu. Il était notre point de contact auprès de Lokad et nous prodiguait une analyse critique, des pistes intéressantes ainsi que les ressources nécessaires à la compréhension des besoins de Lokad. La réunion de mi-parcours avec Joannes Vermorel, CEO de Lokad ainsi que Johannes Lutzeyer, notre référent PSC au département d’informatique nous a permis de redéfinir nos objectifs et d’avoir un premier regard  extérieur au projet.
En plus de ces créneaux, ce projet a également été l’occasion pour nous de nous documenter sur notre temps libre et ensuite implémenter les outils que nous avions conceptualisés durant nos réunions. Grâce à Git, il était assez aisé pour chacun de travailler lorsqu’il le voulait durant la semaine. Nous pouvions ensuite faire le point le mercredi après midi pour échanger sur l’avancement de nos tâches et fixer de nouveaux objectifs pour la semaine d’après.

Le tableau suivant permet de situer l’évolution du PSC tout au long de l’année :

\renewcommand{\arraystretch}{1.5} % Ajuste la hauteur des lignes
\begin{table}[h]
\centering
\begin{tabular}{|l|c|c|c|c|c|c|c|c|}
\hline
\rowcolor{gray!20} \textbf{Tâche} & \textbf{Sept} & \textbf{Oct} & \textbf{Nov} & \textbf{Déc} & \textbf{Janv} & \textbf{Févr} & \textbf{Mars} & \textbf{Avril} \\ \hline
Implémentation outils de base & \cellcolor{gray!50}X & \cellcolor{gray!50}X & & & & & & \\ \hline
Ajout du Grep & & \cellcolor{gray!50}X & & & & & & \\ \hline
Ajout de LangGraph & & & \cellcolor{gray!50}X & \cellcolor{gray!50}X & & & & \\ \hline
RAG Fusion & & & & \cellcolor{gray!50}X & & & & \\ \hline
Pipeline agentique & & & & & \cellcolor{gray!50}X & & & \\ \hline
MàJ chunking parsing & & & & & \cellcolor{gray!50}X & & & \\ \hline
Chunk summarizer & & & & & \cellcolor{gray!50}X & \cellcolor{gray!50}X & & \\ \hline
Reranking & & & & & & \cellcolor{gray!50}X & & \\ \hline
Ajout du Graph tool & & & & & & \cellcolor{gray!50}X & \cellcolor{gray!50}X & \\ \hline
Hyde & & & & & & & \cellcolor{gray!50}X & \\ \hline
Raptor & & & & & & & \cellcolor{gray!50}X & \\ \hline
Benchmarks & & & & & & & & \cellcolor{gray!50}X \\ \hline
\end{tabular}
\caption{Tableau de progression chronologique du PSC}
\end{table}

Nous avons débuté le projet par la conception de notre pipeline de RAG, inspirée de celle du PSC réalisé l’année dernière pour Lokad. Cette vision reposait sur une approche classique et linéaire, telle qu'elle est souvent présentée dans la littérature académique sur le RAG.

Nous nous sommes documentés sur les outils de base qui nous paraissaient les plus adaptés, notamment \textit{FAISS}, \textit{sentence transformers}.
En novembre, nous avions déjà une base de RAG avec un mode verbose et un type de benchmark afin d’observer directement les apports de chaque ajout à venir.

Néanmoins, la confrontation avec les cas d'usage réels a rapidement invalidé cette approche simpliste. La réalité du terrain nous a poussés vers une architecture Agentique non-linéaire.

Une fois celle-ci implémentée, nous avons pu ajouter brique par brique les outils nécessaires à l'amélioration du modèle de manière indépendante. Les efforts sont alors concentrés sur l'optimisation de la performance et la fiabilité des réponses. Il fallait déterminer les configurations les plus efficaces pour chaque outil. Ce perfectionnement pouvait aller de la modification des prompts à l'ajustement de variables comme la taille des chunks lors de la création de l'index.

La répartition des tâches durant tout ce processus s’est effectuée naturellement en fonction des goûts et des capacités de chacun. 

\subsection{Livrables et diffusion}
En complément du manuscrit, nous avons préparé les deux supports publics associés au projet : un dépôt technique, qui regroupe l'implémentation et l'historique de développement, et une page publique plus synthétique, destinée à présenter le sujet, la composition du groupe et l'architecture générale du système.

\begin{itemize}
    \item \textbf{\textcolor{blue!70!black}{Dépôt GitHub du projet}} [\url{https://github.com/ClementLokad/llm-DSL-info-extraction}]
    \item \textbf{\textcolor{green!45!black}{Page publique de présentation}} [\url{https://kpihx.github.io/envision-copilot-presentation/}]
\end{itemize}

Cette séparation nous a permis de distinguer clairement le support de développement complet et le support de communication synthétique demandé pour la restitution finale du PSC.

\subsection{Enseignements tirés}
    Ce projet scientifique collectif nous a appris de nombreuses choses.
Tout d'abord, la partie organisationnelle est bien évidemment au coeur du projet. Avec une équipe plus nombreuse que ce à quoi nous étions habitués dans un cadre scolaire, la bonne coordination n'était pas assurée de prime abord.

Le fait d'appartenir à la même section rendait la communication plus aisée puisque nous nous croisions régulièrement en dehors des réunions. Cela permettait d'échanger régulièrement sur l'avancée du projet et débattre sur les implémentations en cours, poser des questions sur les outils des autres etc.

Le thème du projet a été un des moteurs qui ont conduit à la création du groupe. Nous étions effectivement convaincus que cette problématique particulièrement actuelle nous bénéficierait dans notre avenir professionnel. Nous avons eu l'occasion de découvrir plus en profondeur le fonctionnement des LLMs et l'écosystème en mutation permanente dans lequel ils s'inscrivent. Nous avons également pu renforcer d'autres compétences, maîtrise de Python, de Git.
La complexité de la tâche nous a aussi permis de développer notre persévérance, notre rigueur et notre curiosité scientifique.

En fin de compte, ce travail de longue haleine nous a beaucoup enseigné que ce soit sur le plan social ou scientifique.

\subsection{Conclusion}
Nous sommes positivement surpris par les résultats finaux que nous avons obtenus et apprécions l'intérêt concret que présente le produit à l'état final. Nous espérons que Lokad partage cette satisfaction et que notre outil trouvera une véritable application dans le quotidien des \textit{Supply Chain Scientists} auxquels l'outil était destiné.
Cette satisfaction que nous éprouvons ne signifie pas que l'ensemble de notre parcours fut parfait. Concernant l'organisation, nous ne disposions pas tous du même socle en compétences informatiques. Cela a parfois compliqué la compréhension du code des autres membres, notamment lors de \textit{pushs} abondants. Nous aurions pu éviter ce problème en utilisant plus tôt des outils comme Copilot pour générer de la documentation et commenter le code de manière efficace. L'indépendance des outils a fortement contribué à simplifier les tâches de chacun mais pouvait parfois inciter à ne se focaliser que sur une partie du projet sans bien en connaître le reste.